\section{Discussion}
This section separates observed facts from interpretation scope.

\textbf{Observed in native runs.}
Policy choice changes rank order by workload (Table~\ref{tab:policy_ranking}).
Significant differences are concentrated in two metrics (\texttt{harm\_weighted\_cost}, \texttt{p95\_latency\_ms})
and three workloads (balanced, minority\_spike, skewed\_burst).
Recovery\_stress shows no Holm-significant pairwise differences; the largest observed mean differences are
\(\Delta\texttt{harm\_weighted\_cost}=0.00197\) and \(\Delta\texttt{p95\_latency\_ms}=0.01953\).

\textbf{Observed non-significance that is retained.}
\texttt{drop\_rate} and \texttt{event\_miss\_rate} are identical across non-error policies within each workload.
\texttt{group\_miss\_gap} remains 0.0 for all non-error cells.
\texttt{group\_latency\_gap\_ms} differs numerically but remains non-significant after Holm correction.
These null findings are part of the result, not exclusions.

\textbf{Interpretation boundary.}
The paper reports comparative behavior under this synthetic workload family and fixed metric definitions.
No causal mechanism is claimed beyond observed associations.
No claim is made that one policy is globally superior across all workloads.

\textbf{Limitations.}
Workloads are synthetic, fairness weighting is configuration-dependent, and policy behavior is evaluated in
single-host experimental setups.
Therefore, external validity to different workload distributions or deployment constraints is not asserted.
